\chapter{Abstract}

Medical App V3 is a full-stack medical monitoring graduation project designed to connect patient-facing mobile workflows with ECG and PPG signal acquisition, backend persistence, and AI-assisted physiological analysis. The system addresses a practical problem in remote and semi-continuous health monitoring: physiological readings collected from hardware devices are often difficult to organize, interpret, and present to patients or doctors in a usable software workflow. The implemented project proposes an integrated solution composed of an ESP32-based sensing layer, a FastAPI backend, a PostgreSQL relational database, local Keras model adapters, and a Flutter application for patients, doctors, and administrators.

The hardware firmware samples ECG and PPG values at 250 Hz, groups them into batches, estimates sensor state metadata such as finger detection and PPG availability, and sends JSON payloads to the backend. The backend validates incoming readings with Pydantic schemas, stores sessions and raw readings through SQLAlchemy models, optionally queues analysis jobs, and generates decision-support outputs using deterministic signal processing and optional local AI models. The Flutter frontend presents authentication, profile management, patient and doctor workflows, appointment scheduling, reminders, contacts, device assignment, session listing, ECG/PPG charts, metrics, alerts, and AI summaries.

The project demonstrates an end-to-end medical software architecture that combines mobile usability, IoT data collection, relational persistence, role-aware APIs, and AI-assisted analysis. Its expected impact is to provide a technically coherent prototype for monitoring cardiovascular signals and supporting clinical review. The implementation explicitly states that AI results are decision-support only and not a final medical diagnosis, which is an important limitation and safety boundary for academic and practical evaluation.
